\def\tituloconcepto{Práctica guiada: Bell, variantes y GHZ}
% Si el capitulo se compila solo (sin main.tex), esta cabecera abre el documento.
% Antes de % Si el capitulo se compila solo (sin main.tex), esta cabecera abre el documento.
% Antes de % Si el capitulo se compila solo (sin main.tex), esta cabecera abre el documento.
% Antes de \input{cabecera} cada capitulo define \tituloconcepto.
\ifdefined\enmain\else
  \documentclass[11pt,a4paper]{article}
  \input{preambulo}
  \begin{document}
  \title{\tituloconcepto}
  \author{\autor}
  \date{\today}
  \maketitle
\fi
 cada capitulo define \tituloconcepto.
\ifdefined\enmain\else
  \documentclass[11pt,a4paper]{article}
  % Preambulo compartido por todos los documentos. Compilar con XeLaTeX/LuaLaTeX
% (o Tectonic). Si se usa Overleaf, elegir el compilador XeLaTeX.
\usepackage[spanish,es-noshorthands]{babel}
\usepackage{amsmath,amssymb}
\usepackage{booktabs,array}
\usepackage[a4paper,margin=2.5cm]{geometry}
\usepackage{xcolor}
\usepackage{listings}
\usepackage[most]{tcolorbox}
\usepackage{quantikz}
\usepackage[hidelinks]{hyperref}

% ---- Autor (editar a gusto) -------------------------------------------------
\newcommand{\autor}{Genesis Silesky}

% ---- Notacion de Dirac y atajos ---------------------------------------------
% quantikz ya define \ket y \bra; DeclareRobustCommand los redefine sin error.
\DeclareRobustCommand{\ket}[1]{\left|#1\right\rangle}
\DeclareRobustCommand{\bra}[1]{\left\langle#1\right|}
\DeclareRobustCommand{\braket}[2]{\left\langle#1\middle|#2\right\rangle}
\newcommand{\C}{\mathbb{C}}
\newcommand{\R}{\mathbb{R}}
\DeclareMathOperator{\Tr}{Tr}

% ---- Cajas de apoyo --------------------------------------------------------
% "concepto": definiciones previas que hacen falta para entender el capitulo.
\newtcolorbox{concepto}[1]{breakable,colback=blue!4,colframe=blue!55!black,
  fonttitle=\bfseries,title={Concepto previo: #1}}
% "nota": aclaraciones y advertencias.
\newtcolorbox{nota}[1]{breakable,colback=orange!6,colframe=orange!70!black,
  fonttitle=\bfseries,title={#1}}

% ---- Estilos de listados ----------------------------------------------------
\lstdefinestyle{py}{language=Python,basicstyle=\ttfamily\footnotesize,
  breaklines=true,postbreak=\mbox{\textcolor{gray}{$\hookrightarrow$}\space},
  numbers=left,numberstyle=\tiny\color{gray},frame=single,
  commentstyle=\color{green!40!black}\itshape,
  keywordstyle=\color{blue!70!black}\bfseries,
  stringstyle=\color{red!60!black},showstringspaces=false,tabsize=4,
  columns=fullflexible,keepspaces=true}
\lstdefinestyle{salida}{basicstyle=\ttfamily\footnotesize,breaklines=true,
  frame=leftline,backgroundcolor=\color{gray!8},columns=fullflexible,
  keepspaces=true}

  \begin{document}
  \title{\tituloconcepto}
  \author{\autor}
  \date{\today}
  \maketitle
\fi
 cada capitulo define \tituloconcepto.
\ifdefined\enmain\else
  \documentclass[11pt,a4paper]{article}
  % Preambulo compartido por todos los documentos. Compilar con XeLaTeX/LuaLaTeX
% (o Tectonic). Si se usa Overleaf, elegir el compilador XeLaTeX.
\usepackage[spanish,es-noshorthands]{babel}
\usepackage{amsmath,amssymb}
\usepackage{booktabs,array}
\usepackage[a4paper,margin=2.5cm]{geometry}
\usepackage{xcolor}
\usepackage{listings}
\usepackage[most]{tcolorbox}
\usepackage{quantikz}
\usepackage[hidelinks]{hyperref}

% ---- Autor (editar a gusto) -------------------------------------------------
\newcommand{\autor}{Genesis Silesky}

% ---- Notacion de Dirac y atajos ---------------------------------------------
% quantikz ya define \ket y \bra; DeclareRobustCommand los redefine sin error.
\DeclareRobustCommand{\ket}[1]{\left|#1\right\rangle}
\DeclareRobustCommand{\bra}[1]{\left\langle#1\right|}
\DeclareRobustCommand{\braket}[2]{\left\langle#1\middle|#2\right\rangle}
\newcommand{\C}{\mathbb{C}}
\newcommand{\R}{\mathbb{R}}
\DeclareMathOperator{\Tr}{Tr}

% ---- Cajas de apoyo --------------------------------------------------------
% "concepto": definiciones previas que hacen falta para entender el capitulo.
\newtcolorbox{concepto}[1]{breakable,colback=blue!4,colframe=blue!55!black,
  fonttitle=\bfseries,title={Concepto previo: #1}}
% "nota": aclaraciones y advertencias.
\newtcolorbox{nota}[1]{breakable,colback=orange!6,colframe=orange!70!black,
  fonttitle=\bfseries,title={#1}}

% ---- Estilos de listados ----------------------------------------------------
\lstdefinestyle{py}{language=Python,basicstyle=\ttfamily\footnotesize,
  breaklines=true,postbreak=\mbox{\textcolor{gray}{$\hookrightarrow$}\space},
  numbers=left,numberstyle=\tiny\color{gray},frame=single,
  commentstyle=\color{green!40!black}\itshape,
  keywordstyle=\color{blue!70!black}\bfseries,
  stringstyle=\color{red!60!black},showstringspaces=false,tabsize=4,
  columns=fullflexible,keepspaces=true}
\lstdefinestyle{salida}{basicstyle=\ttfamily\footnotesize,breaklines=true,
  frame=leftline,backgroundcolor=\color{gray!8},columns=fullflexible,
  keepspaces=true}

  \begin{document}
  \title{\tituloconcepto}
  \author{\autor}
  \date{\today}
  \maketitle
\fi


\section{Práctica guiada: Bell, variantes y GHZ}

\subsection{Objetivo}
Ejecutar el circuito de Bell, predecir resultados antes de verlos, modificarlo
para distinguir superposición de determinismo, y generalizarlo a tres qubits
(estado GHZ).

\subsection{Conceptos previos}
\begin{concepto}{entrelazamiento y shots}
Haber leído el capítulo~05b (estado de Bell) y el 04 (shots y frecuencia).
\end{concepto}

\subsection{Fundamento teórico}
\begin{itemize}
  \item \textbf{Bell:} $H$ en $q_0$ y CNOT $q_0\to q_1$ dan
    $\tfrac1{\sqrt2}(\ket{00}+\ket{11})$. Se esperan $\approx50\%$ de $00$ y
    $\approx50\%$ de $11$.
  \item \textbf{Variante con $X$:} $X\ket0=\ket1$ no crea superposición; el
    CNOT copia ese $1$, así que el estado es $\ket{11}$ y el resultado es
    determinista. Esto aísla el papel de $H$: es la fuente de la aleatoriedad.
  \item \textbf{GHZ:} con $H$ y dos CNOT en cadena se obtiene
    $\tfrac1{\sqrt2}(\ket{000}+\ket{111})$.
\end{itemize}
\begin{center}
\begin{quantikz}
\lstick{$q_0$} & \gate{H} & \ctrl{1} & \qw      & \meter{}\\
\lstick{$q_1$} & \qw      & \targ{}  & \ctrl{1} & \meter{}\\
\lstick{$q_2$} & \qw      & \qw      & \targ{}  & \meter{}
\end{quantikz}
\end{center}

\subsection{Código (completamente comentado)}
\lstinputlisting[style=py]{../07_bell_practica.py}

\subsection{Explicación celda por celda}
\begin{description}
  \item[Celda 2] Bell con mediciones; \texttt{measure([0,1],[0,1])} conecta
    cada qubit con un bit clásico.
  \item[Celda 3] Sustituye \texttt{h(0)} por \texttt{x(0)}.
  \item[Celda 4] Reto: GHZ de 3 qubits.
  \item[Celda 5] Reto: distintos números de shots.
\end{description}

\subsection{Salida obtenida}
\lstinputlisting[style=salida]{../salidas/07_bell_practica.txt}
(El dibujo de texto del circuito, que imprime el script, no se reproduce aquí.)

\subsection{Interpretación}
\begin{itemize}
  \item Con $x$ en lugar de $h$ solo sale $11$: ya no hay aleatoriedad.
  \item GHZ solo da $000$ y $111$.
  \item Con $10$ shots, las proporciones pueden alejarse mucho de $50/50$
    (aquí, 8 contra 2); con $10\,000$ se acercan. Es ruido \emph{estadístico},
    no físico (el siguiente capítulo introduce el ruido físico).
\end{itemize}
\textbf{Pregunta clave:} ¿por qué nunca sale $01$ ni $10$? Porque el CNOT
fuerza que $q_1$ copie el valor de $q_0$ en la base computacional.

\subsection{Ejercicios}
\begin{enumerate}
  \item Prediga y verifique qué ocurre con $h(0)$ seguido de $x(1)$ (sin CNOT).
  \item Extienda GHZ a 5 qubits.
  \item Mida solo uno de los tres qubits de GHZ y observe cuál es la
    distribución del resto.
\end{enumerate}

\ifdefined\enmain\else
  \end{document}
\fi

